\chapter{Spinal Fusion}

\section*{Introduction \& Clinical Indications}
Spinal fusion is a complex surgical procedure aimed at permanently joining two or more vertebrae in the spine to eliminate motion between them. This stabilization is achieved through a biological process where bone grafts are utilized to promote the formation of a solid bony bridge, known as arthrodesis. The procedure is performed in various anatomical regions, including the cervical, thoracic, and lumbar spine, depending on the underlying pathology. Spinal fusion is clinically significant as it addresses conditions that lead to spinal instability, deformity, or debilitating pain, often resulting from degenerative diseases, trauma, or tumors. The procedure may involve the use of internal fixation devices, such as screws, rods, and interbody cages, which provide immediate stability while the bone graft heals and fuses.

\begin{itemize}
  \item Spondylodesis
  \item Arthrodesis of the spine
  \item Posterior Lumbar Interbody Fusion (PLIF)
  \item Transforaminal Lumbar Interbody Fusion (TLIF)
  \item Anterior Lumbar Interbody Fusion (ALIF)
  \item Extreme Lateral Interbody Fusion (XLIF)
  \item Posterolateral Fusion (PLF)
\end{itemize}

The primary principle of spinal fusion is to create a stable environment that promotes bone growth across adjacent vertebral segments, resulting in a single, rigid unit. This process, known as osteointegration, involves the placement of bone graft material, which serves as a scaffold for new bone formation. The graft can be autograft (from the patient's own body), allograft (donated human bone), or synthetic substitutes. 

Instrumentation, including pedicle screws, rods, and interbody cages, is utilized to provide immediate mechanical stability, holding the vertebrae in a corrected position and protecting the graft during the healing phase. The surgical rationale is to eliminate painful motion, correct deformities, or stabilize segments that are unstable due to trauma, degeneration, or tumors. By fusing the vertebrae, the source of pain originating from the motion segment, such as degenerative disc disease or spondylolisthesis, is theoretically removed. The technical goals include achieving solid bony union, restoring spinal alignment, decompressing neural elements if necessary, and maintaining neurological function.

The concept of spinal fusion has evolved significantly since its inception in the early 20th century. In 1911, Russell Hibbs and Fred Albee independently performed the first successful spinal fusions to treat Pott's disease (spinal tuberculosis). Hibbs' technique involved fusing the posterior elements, while Albee utilized a tibial bone graft placed into split spinous processes. These pioneering efforts marked a significant advancement in the treatment of spinal instability and deformity.

In the 1950s, Paul Harrington introduced Harrington rods for scoliosis correction, providing significant distraction and stabilization. The development of segmental instrumentation in the 1960s and 1970s allowed for more rigid fixation and better deformity correction. Pioneers like Roy-Camille and Cotrel-Dubousset revolutionized spinal surgery with pedicle screw fixation, enabling robust stabilization and correction of complex deformities. The late 20th and early 21st centuries have seen further advancements, including interbody cages, minimally invasive techniques, and biological bone graft substitutes, aimed at improving fusion rates and reducing surgical morbidity.

Spinal fusion is indicated for various conditions that cause spinal instability, deformity, or intractable pain unresponsive to conservative management. Common clinical indications include:

\begin{itemize}
  \item Stabilization of the spine in degenerative spondylolisthesis, where one vertebra slips over another, causing pain and nerve compression.
  \item Correction of severe spinal deformities such as scoliosis or kyphosis that are progressive or symptomatic.
  \item Treatment of spinal instability due to trauma, such as fractures or dislocations, compromising vertebral integrity.
  \item Alleviation of chronic back or neck pain from severe degenerative disc disease, particularly with instability or nerve impingement.
  \item Stabilization after the removal of a spinal tumor or infection that has compromised vertebral integrity.
  \item Addressing pseudoarthrosis, where a previous spinal fusion has failed to achieve solid bony union.
  \item Decompression and stabilization in cases of severe spinal stenosis with associated instability.
  \item Treatment of congenital spinal anomalies leading to progressive deformity or neurological deficits.
\end{itemize}

Spinal fusion can serve as both a primary and secondary surgical intervention, depending on the underlying pathology. For progressive spinal deformities, such as severe scoliosis or acute spinal trauma with instability, it is often considered a primary surgical treatment. In these cases, the immediate need for stabilization and correction outweighs prolonged conservative measures. Conversely, for degenerative conditions like chronic low back pain due to degenerative disc disease, spinal fusion is typically a secondary or salvage procedure, reserved for patients who have failed extensive non-operative treatments and whose symptoms significantly impair their quality of life.

\section*{Relevant Surgical Anatomy}
The anatomy relevant to spinal fusion encompasses the vertebrae, intervertebral discs, spinal cord, nerve roots, and surrounding soft tissues. The vertebral column consists of 33 vertebrae, categorized into cervical, thoracic, lumbar, sacral, and coccygeal regions. Each vertebra is composed of a vertebral body, pedicles, laminae, spinous processes, and transverse processes. 

The intervertebral discs, situated between adjacent vertebrae, consist of an outer annulus fibrosus and an inner nucleus pulposus, providing cushioning and flexibility. The arterial supply to the spine is primarily through the segmental arteries, which branch from the aorta, while venous drainage occurs via the internal and external vertebral venous plexuses. 

Nerve relations are critical, as the spinal cord runs within the vertebral canal, and spinal nerves exit through the intervertebral foramina. Key surgical landmarks include McBurney's point for lumbar approaches and Calot's triangle for cervical approaches. Understanding these anatomical structures is essential for minimizing complications during surgery.

\begin{itemize}
  \item Vertebrae: Composed of body, pedicles, laminae, and processes.
  \item Intervertebral discs: Annulus fibrosus and nucleus pulposus.
  \item Arterial supply: Segmental arteries from the aorta.
  \item Venous drainage: Internal and external vertebral venous plexuses.
  \item Nerve relations: Spinal cord and exiting spinal nerves.
\end{itemize}

\section*{Preoperative Workup \& Preparation}
Thorough preoperative preparation is essential to optimize outcomes and minimize risks. Key components include:

\begin{itemize}
  \item \textbf{Medical Evaluation:} Comprehensive workup including blood tests, ECG, and chest X-ray; cardiac and pulmonary clearance may be required.
  \item \textbf{Medication Review:} Patients should stop blood thinners and herbal supplements 5-7 days prior to surgery.
  \item \textbf{Smoking Cessation:} Ideally, patients should cease smoking several months before surgery to enhance healing.
  \item \textbf{NPO Status:} Patients must be NPO for 6-8 hours before surgery to prevent aspiration.
  \item \textbf{Bowel Preparation:} Required for anterior approaches to reduce infection risk.
  \item \textbf{Preoperative Antibiotics:} Administered shortly before the incision to minimize infection risk.
  \item \textbf{Blood Typing and Crossmatch:} Prepared in case a transfusion is needed.
  \item \textbf{Informed Consent:} Detailed discussion regarding the procedure, risks, and alternatives.
\end{itemize}

\textbf{Absolute Contraindications:}
\begin{itemize}
  \item Active systemic or local spinal infection, significantly increasing surgical site infection risk.
  \item Severe, uncontrolled coagulopathy that cannot be managed preoperatively.
  \item Uncontrolled systemic medical conditions that pose high surgical risks.
  \item Severe osteoporosis preventing stable fixation of instrumentation.
  \item Malignancy with a very short life expectancy where surgery risks outweigh benefits.
\end{itemize}

\textbf{Relative Contraindications and Precautions:}
\begin{itemize}
  \item Morbid obesity, increasing surgical complexity and complication rates.
  \item Active substance abuse or uncontrolled psychiatric illness complicating recovery.
  \item Severe peripheral vascular disease, especially for anterior approaches.
  \item Significant nutritional deficiencies impairing healing.
  \item Previous spinal surgery at the same level, increasing surgical difficulty.
  \item Advanced age with multiple comorbidities requiring careful assessment.
\end{itemize}

\section*{Surgical Technique \& Procedure}
Spinal fusion is performed under general anesthesia, ensuring the patient is unconscious and pain-free. The surgical approach varies based on the targeted spinal region: prone for posterior approaches (PLIF, TLIF, PLF), supine for anterior approaches (ALIF), or lateral decubitus for lateral approaches (XLIF, LLIF). Intraoperative imaging, such as fluoroscopy, is routinely used to confirm spinal levels and guide instrumentation placement.

The general steps for a posterior lumbar interbody fusion (PLIF) or transforaminal lumbar interbody fusion (TLIF) typically include:

\begin{itemize}
  \item \textbf{Incision and Exposure:} A midline incision is made, and the paraspinal muscles are dissected and retracted to expose the posterior elements of the vertebrae.
  \item \textbf{Decompression (if needed):} A laminectomy, facetectomy, or foraminotomy may be performed to relieve pressure on the spinal cord or nerve roots.
  \item \textbf{Discectomy and Endplate Preparation:} The intervertebral disc is removed, and the cartilaginous endplates are prepared to expose bleeding bone.
  \item \textbf{Bone Graft and Interbody Cage Placement:} An interbody cage packed with bone graft material is inserted into the disc space to restore height and provide support.
  \item \textbf{Pedicle Screw Insertion:} Pedicle screws are placed into the vertebral pedicles above and below the fused segment.
  \item \textbf{Rod Placement and Compression/Distraction:} Metal rods are connected to the screws, and the construct is adjusted for optimal alignment.
  \item \textbf{Posterolateral Grafting (optional):} Additional bone graft may be placed along the transverse processes to enhance fusion.
  \item \textbf{Wound Closure:} The incision is closed in layers, and a drain may be placed to collect excess fluid.
\end{itemize}

The specific details vary based on the number of levels fused, the approach, and the patient's anatomy.

\section*{Complications \& Risks}
Spinal fusion carries potential risks, categorized into general surgical complications and procedure-specific risks.

\begin{itemize}
  \item \textbf{General Surgical Risks:}
    \begin{itemize}
      \item \textbf{Bleeding:} Significant blood loss requiring transfusion.
      \item \textbf{Infection:} Surgical site infections that may necessitate further intervention.
      \item \textbf{Anesthesia Complications:} Adverse reactions to anesthesia.
      \item \textbf{Deep Vein Thrombosis (DVT) and Pulmonary Embolism (PE):} Potentially life-threatening blood clots.
      \item \textbf{Urinary Tract Infection (UTI) or Pneumonia:} Common postoperative infections.
    \end{itemize}
  
  \item \textbf{Procedure-Specific Risks:}
    \begin{itemize}
      \item \textbf{Pseudarthrosis (Failure of Fusion):} Common long-term complication requiring revision.
      \item \textbf{Hardware Failure:} Breakage or loosening of screws or rods.
      \item \textbf{Nerve Root or Spinal Cord Injury:} Potential new or worsened neurological deficits.
      \item \textbf{Dural Tear (Cerebrospinal Fluid Leak):} Leading to headaches or meningitis.
      \item \textbf{Adjacent Segment Disease (ASD):} Degeneration at segments adjacent to the fusion.
      \item \textbf{Chronic Pain:} Persistent pain despite successful fusion.
      \item \textbf{Graft Site Pain:} Chronic pain at the donor site if autograft is used.
      \item \textbf{Vascular Injury:} Risk of injury to major blood vessels during anterior approaches.
      \item \textbf{Retrograde Ejaculation:} Specific risk for males undergoing anterior lumbar fusion.
      \item \textbf{Vision Loss:} Rare but serious complication from prolonged prone positioning.
    \end{itemize}
\end{itemize}

\section*{Postoperative Care \& Recovery}
Postoperative recovery following spinal fusion involves several phases. Immediately after surgery, patients are transferred to the Post-Anesthesia Care Unit (PACU) for monitoring. Vital signs, neurological status, and pain levels are assessed continuously. Pain management is crucial, often utilizing a multimodal approach with intravenous opioids and non-opioid analgesics.

During the first 24-48 hours, the focus is on pain control, early mobilization, and preventing complications. Patients are encouraged to ambulate with assistance on the first postoperative day, which helps prevent complications such as DVT and pneumonia. Physical therapists instruct patients on proper body mechanics to avoid twisting or bending the spine. 

Patients are typically discharged within 3-7 days, depending on the extent of surgery and recovery progress. Long-term recovery involves gradually increasing activity levels while avoiding bending, twisting, and heavy lifting for several months to allow the fusion to solidify. Formal physical therapy usually begins a few weeks post-surgery to strengthen core muscles and improve flexibility. Full recovery and solid fusion can take 6-12 months or longer.

During spinal fusion surgery, the surgeon documents the condition of the vertebrae, intervertebral discs, and any associated pathologies. The pathology report on resected tissues may reveal degenerative changes, tumors, or signs of infection. Successful fusion is characterized by the formation of a solid bony bridge between the fused vertebrae, with the absence of motion on dynamic imaging and evidence of continuous bony bridging on radiological assessments.

A normal postoperative course following spinal fusion is characterized by a gradual reduction in pain, improvement in neurological symptoms, and enhanced functional capacity. Patients typically experience significant pain relief within weeks to months, allowing for a return to daily activities. 

The timeline for recovery varies based on individual factors, including the extent of surgery and the patient's overall health. Most patients can expect to resume light activities within a few weeks, with more strenuous activities gradually reintroduced over several months. Regular follow-up appointments are essential to monitor healing and assess fusion progress.

Postoperative care is crucial for ensuring optimal recovery. Follow-up visits are typically scheduled at 2 weeks, 6 weeks, 3 months, 6 months, and 1 year post-surgery. Key components of postoperative care include:

\begin{itemize}
  \item \textbf{Surgical Check-ups:} Regular visits to assess recovery and fusion progress.
  \item \textbf{Imaging Studies:} X-rays are taken at follow-up visits to evaluate fusion and instrumentation integrity.
  \item \textbf{Suture/Staple Removal:} Usually performed at the 2-week follow-up visit.
  \item \textbf{Physical Therapy and Rehabilitation:} Initiated a few weeks post-surgery to strengthen muscles and improve mobility.
  \item \textbf{Activity Resumption Guidelines:} Specific instructions on gradually resuming normal activities.
  \item \textbf{Pain Management:} Ongoing management of postoperative pain.
  \item \textbf{Warning Signs:} Patients are educated on signs to report, such as fever, worsening pain, or neurological changes.
\end{itemize}

\section*{Outcomes \& Prognosis}
Spinal fusion is one of the most common spinal surgical procedures performed globally, with its incidence steadily increasing, particularly among older adults. The highest rates are observed in the lumbar spine, primarily for degenerative conditions such as degenerative disc disease, spondylolisthesis, and spinal stenosis. 

The procedure predominantly affects individuals over 50 years of age, with a slight female predominance for degenerative spondylolisthesis. Cervical fusion is commonly performed for degenerative disc disease and myelopathy. While scoliosis correction is less frequent, it remains a significant indication for fusion in adolescents and young adults. The overall mortality rate associated with elective spinal fusion is low, typically less than 0.5\%, but morbidity rates can be substantial, ranging from 10\% to 20\%, encompassing complications such as infection and pseudarthrosis. The increasing prevalence is attributed to an aging population, improved diagnostic capabilities, and evolving surgical techniques.

The outcomes and prognosis following spinal fusion are generally favorable for appropriately selected patients. Success rates for significant pain relief and functional improvement typically range from 70\% to 90\%. Patients undergoing fusion for spinal deformity correction often experience substantial improvement in alignment and pain reduction. 

Factors influencing prognosis include smoking status, obesity, the number of levels fused, comorbidities, and preoperative neurological status. While fusion aims for permanent stability, long-term complications such as adjacent segment disease can occur, with an incidence of 2-4\% per year. Pseudarthrosis rates vary but can be as high as 10-20\%, particularly in smokers or multi-level fusions. Despite these risks, the majority of patients experience improved quality of life and functional independence, with a return to many pre-surgical activities.

\section*{References}
\begin{enumerate}
  \item MedlinePlus. Spinal Fusion. U.S. National Library of Medicine; 2024. Available from: \url{https://medlineplus.gov/spinalfusion.html}
  \item Chapman JR, et al. Spinal Fusion. In: Orthopaedic Knowledge Update: Spine 6. American Academy of Orthopaedic Surgeons; 2020.
  \item Vaccaro AR, et al. Rothman-Simeone and Herkowitz's The Spine. 8th ed. Elsevier; 2023.
  \item North American Spine Society (NASS). Clinical Guidelines for the Diagnosis and Treatment of Degenerative Lumbar Spondylolisthesis. NASS; 2021.
\end{enumerate}

\clearpage
